\documentclass[11pt]{article}
\usepackage{../homework}
\profname{Dr. Claire Zukowski}
\classcode{PHYS 331 Quantum Mechanics I}
\assignnum{Problem Set 6}
\duedate{4 November 2026}
\begin{document}

\begin{problem}{1}
  Calculate \(\ev{x}\), \(\ev{p}\), \(\ev{x^2}\), \(\ev{p^2}\), \(\Delta x\), and \(\Delta p\) for the \(n^{\text{th}}\) excited energy eigenstate of the infinite square well.
  Check that the uncertainty principle is satisfied.
  What state comes closest to saturating the uncertainty limit \(\Delta x\Delta p \geq \frac{\hbar}{2}\)?
\end{problem}
\begin{answer}
  The \( n^{th} \) eigenstate for the infinite square well for a particule of mass \( n \) in a well of length \( L \) is
  \begin{equation}
  \label{eq:1}
  \varphi_n = \sqrt{\frac{2}{L}}\sin(\frac{n\pi x}{L}), \quad E_n = \frac{n^2 \pi^2 \hbar^2}{2mL^2}.
\end{equation}

We then calculate,
\begin{equation}
\label{eq:2}
\begin{aligned}
  \ev{x} &= \ev{x}{\varphi} = \frac{2}{L} \int\limits_0^L x\sin^2\left(\frac{n\pi x}{L}\right) \dd{x} = \frac{2}{L}\left( \frac{L^2}{4} \right) \\
  &= \frac{L}{2}.
\end{aligned}
\end{equation}

\begin{equation}
\label{eq:3}
\begin{aligned}
  \ev{p} &= \ev{p}{\varphi} = \frac{2}{L} \int\limits_0^L \sin(\frac{n\pi x}{L}) i\hbar\frac{n\pi}{L}\cos(\frac{n\pi x}{L})\dd{x} = 0,
\end{aligned}
\end{equation}
where the last equality follows from the orthogonality of sine and cosine.

\begin{equation}
\label{eq:4}
\begin{aligned}
  \ev{x^2} &= \ev{x^2}{\varphi} = \frac{2}{L} \int\limits_0^L x^2 \sin^2(\frac{n\pi x}{L}) \dd{x} \\
           &= \frac{2}{L}\left( L^3(\frac{1}{6} - \frac{1}{4\pi^2n^2}) \right) = L^2\left(\frac{1}{3} - \frac{1}{2\pi^2 n^2}\right).
\end{aligned}
\end{equation}
\begin{equation}
\label{eq:5}
\begin{aligned}
  \ev{p^2} &= \ev{p^2}{\varphi} = \frac{2}{L} \int\limits_0^L \sin(\frac{n\pi x}{L})\frac{n^2\hbar^2\pi^2}{L^2}\sin(\frac{n\pi x}{L}) \dd{x} \\
  &= \frac{n^2 \pi^2 \hbar^2}{L^2}.
\end{aligned}
\end{equation}

With these values we can calculate
\begin{equation}
\label{eq:6}
\Delta x = \sqrt{\ev{x^2} - \ev{x}^2} = L\sqrt{\frac{1}{3} - \frac{1}{2\pi^2n^2} - \frac{1}{4}}
\end{equation}
\begin{equation}
\label{eq:7}
\Delta p = \sqrt{\ev{p^2} - \ev{p}^2}
= \frac{n\pi\hbar}{L}.
\end{equation}

Thus,
\begin{equation}
\label{eq:8}
\begin{aligned}
\Delta x \Delta p
&= L\sqrt{\frac{1}{12} - \frac{1}{2\pi^2n^2}}
\frac{n\pi\hbar}{L} \\
&= \hbar\sqrt{\frac{n^2\pi^2}{12} - \frac{1}{2}}
\geq \frac{\hbar}{2},
\end{aligned}
\end{equation}
where the inequality follows since \(n \geq 1\) and
\(\pi^2 > 9\).

Finally, since \(\Delta x\Delta p\) is increasing in
\(n\), the ground state \(n=1\) comes closest to
saturating the uncertainty principle, with
\begin{equation}
\label{eq:9}
\Delta x\Delta p
= \hbar\sqrt{\frac{\pi^2}{12} - \frac{1}{2}}.
\end{equation}
\end{answer}

\begin{problem}{2}
  A particle of mass \(m\) is in an infinite square well (\(V=0\) from \(x=0\) to \(a\) and \(V=\infty\) elsewhere).
  You may use all results from class/textbook without re-deriving them.
  Suppose the particle has been prepared in a (very roughly) localized state
  \begin{equation}
    \label{eq:p2-initial-wavefunction}
    \psi(x,t=0)=
    \begin{cases}
      A\,x/a & 0<x<a,\\
      0 & \text{elsewhere}.
    \end{cases}
  \end{equation}
  \begin{enumerate}[label=(\alph*)]
    \item Determine the normalization constant, \(A\).
\begin{answer}
    We normalize the initial wave function by requiring
    \begin{equation}
      \label{eq:p2-normalization}
      \begin{aligned}
        1
        &= \int_{-\infty}^{\infty}
        |\psi(x,0)|^2\dd{x} \\
        &= \frac{|A|^2}{a^2}\int_0^a x^2\dd{x} \\
        &= \frac{|A|^2a}{3}.
      \end{aligned}
    \end{equation}
    Choosing \(A\) real and positive gives
    \begin{equation}
      \label{eq:p2-A}
      \boxed{A=\sqrt{\frac{3}{a}}}.
    \end{equation}
          \end{answer}
    \item The initial state can be expanded in the eigenstates of the infinite well, \(\psi(x,t=0)=\sum_n c_n\phi_n(x)\).
      Derive an expression, in terms of \(n\), for the \(c_n\) coefficients.
          \begin{answer}
   Recall that the normalized eigenfunctions of the
    infinite square well are
    \begin{equation}
      \label{eq:p2-eigenfunctions}
      \phi_n(x)=\sqrt{\frac{2}{a}}
      \sin\left(\frac{n\pi x}{a}\right).
    \end{equation}
    By orthonormality, the expansion coefficients are
    \begin{equation}
      \label{eq:p2-cn}
      \begin{aligned}
        c_n
        &= \int_0^a \phi_n^*(x)\psi(x,0)\dd{x} \\
        &= \frac{\sqrt{6}}{a^2}
        \int_0^a x\sin\left(\frac{n\pi x}{a}\right)
        \dd{x}.
      \end{aligned}
    \end{equation}

    Integrating by parts,
    \begin{equation}
      \label{eq:p2-integral}
      \begin{aligned}
        &\int_0^a x\sin\left(\frac{n\pi x}{a}\right)
        \dd{x} \\
        &\quad =
        \left[
        -\frac{ax}{n\pi}
        \cos\left(\frac{n\pi x}{a}\right)
        +\frac{a^2}{n^2\pi^2}
        \sin\left(\frac{n\pi x}{a}\right)
        \right]_0^a \\
        &\quad =-\frac{a^2}{n\pi}(-1)^n.
      \end{aligned}
    \end{equation}

    Substituting gives
    \begin{equation}
      \label{eq:p2-coefficients}
      \boxed{c_n=\frac{(-1)^{n+1}\sqrt{6}}{n\pi}}.
    \end{equation}
          \end{answer}

    \item Compute the probability for measuring the energy to be \(9\pi^2\hbar^2/(2ma^2)\) at \(t=0\).
          \begin{answer}
  The energy eigenvalues are
    \begin{equation}
      \label{eq:p2-energies}
      E_n=\frac{n^2\pi^2\hbar^2}{2ma^2}.
    \end{equation}
    Thus, the specified energy corresponds to \(n=3\).

    The probability of measuring this energy is
    \begin{equation}
      \label{eq:p2-probability}
      \begin{aligned}
        P(E_3)
        &= |c_3|^2 \\
        &= \left|
        \frac{(-1)^4\sqrt{6}}{3\pi}
        \right|^2 \\
        &= \boxed{\frac{2}{3\pi^2}}.
      \end{aligned}
    \end{equation}
          \end{answer}

    \item As an approximation, let's truncate the expansion at \(n=3\) (including the \(n=1,2,3\) terms).
      Write an expression for the full, time-dependent wave function.
          \begin{answer}
    The time evolution of an energy eigenstate is
    \begin{equation}
      \label{eq:p2-time-evolution}
      \phi_n(x,t)=\phi_n(x)e^{-iE_nt/\hbar}.
    \end{equation}
    Therefore, the full wave function is
    \begin{equation}
      \label{eq:p2-full-series}
      \psi(x,t)
      =\sum_{n=1}^{\infty}
      c_n\phi_n(x)e^{-iE_nt/\hbar}.
    \end{equation}

    Truncating at \(n=3\) gives
    \begin{equation}
      \label{eq:p2-truncated}
      \begin{aligned}
        \psi(x,t)\approx{}&
        c_1\phi_1(x)e^{-iE_1t/\hbar}
        +c_2\phi_2(x)e^{-iE_2t/\hbar}\\
        &+c_3\phi_3(x)e^{-iE_3t/\hbar}.
      \end{aligned}
    \end{equation}

    Substituting the coefficients and eigenfunctions,
    \begin{equation}
      \label{eq:p2-final}
      \boxed{
      \begin{aligned}
        \psi(x,t)\approx
        \frac{\sqrt{12/a}}{\pi}
        \bigg[
        &\sin\left(\frac{\pi x}{a}\right)
        e^{-i\pi^2\hbar t/(2ma^2)}\\
        &-\frac12\sin\left(\frac{2\pi x}{a}\right)
        e^{-i4\pi^2\hbar t/(2ma^2)}\\
        &+\frac13\sin\left(\frac{3\pi x}{a}\right)
        e^{-i9\pi^2\hbar t/(2ma^2)}
        \bigg],
      \end{aligned}
      }
    \end{equation}
    for \(0<x<a\), and zero elsewhere.
          \end{answer}
  \end{enumerate}
\end{problem}

\begin{problem}{3}
  Consider a ``half-infinite'' square well in which the potential energy is \(\infty\) for \(x<0\), zero for \(0<x<a\), and finite (with a value \(V\)) for \(x>a\).
  \begin{enumerate}[label=(\alph*)]
    \item The energy eigenstates are almost exactly the same as one set of the solutions to the finite well problem (i.e., either McIntyre Eq.~5.80 or 5.81, with corresponding eigenvalue Eq.~5.83 or 5.85).
      WHICH of those equations are the right ones here?
      Briefly, explain.
      There is going to be one minor difference from the McIntyre solution (in normalization).
      What is that difference?
\begin{answer}
    The correct solutions are the odd eigenfunctions
    in McIntyre Eq.~5.81, with the corresponding
    energy eigenvalue condition in Eq.~5.85.

    Since the potential is infinite for \(x<0\),
    the wave function must vanish at \(x=0\).
    Of the two families of solutions for the
    symmetric finite square well, only the odd
    solutions satisfy this condition, since
    \(\sin(0)=0\), whereas \(\cos(0)=1\).

    The bound-state wave function therefore has
    the form
    \begin{equation}
      \label{eq:p3-wavefunction}
      \psi(x)=
      \begin{cases}
        0, & x\leq 0,\\
        C\sin(kx), & 0<x<a,\\
        C\sin(ka)e^{-q(x-a)}, & x\geq a,
      \end{cases}
    \end{equation}
    where
    \begin{equation}
      \label{eq:p3-parameters}
      k=\frac{\sqrt{2mE}}{\hbar},
      \qquad
      q=\frac{\sqrt{2m(V-E)}}{\hbar}.
    \end{equation}

    Continuity of the derivative at \(x=a\) gives
    \begin{equation}
      \label{eq:p3-boundary}
      kC\cos(ka)=-qC\sin(ka),
    \end{equation}
    and hence
    \begin{equation}
      \label{eq:p3-energy-condition}
      \boxed{-k\cot(ka)=q},
    \end{equation}
    precisely the condition in McIntyre Eq.~5.85.

    The only difference in normalization is that
    the original odd eigenfunctions are normalized
    over the entire real line, whereas here the
    wave function is supported only on \(x>0\).

    Since the probability density of an odd
    eigenfunction is even, its integral over
    \(x>0\) is exactly one-half of its integral
    over the entire real line. Therefore, the
    normalization constant must be larger by
    a factor of \(\sqrt{2}\).

    Equivalently, if \(\varphi_{\mathrm{odd}}\)
    denotes the normalized odd finite-well
    eigenfunction from McIntyre, then
    \begin{equation}
      \label{eq:p3-normalization}
      \boxed{
      \psi(x)=
      \begin{cases}
        0, & x\leq 0,\\
        \sqrt{2}\,\varphi_{\mathrm{odd}}(x),
        & x>0.
      \end{cases}
      }
    \end{equation}
\end{answer}
    \item Give a qualitative (graphical) argument that the energy eigenvalues of this finite well are close to the usual results for an infinite well where \(V_0\) goes to \(\infty\) at \(x=0\) and \(x=a\).
      \begin{itemize}[label=--]
        \item Are the finite well energies a little smaller or a little larger than the ``ideal'' infinite well energies?
          Why?
        \item As you go higher and higher in (allowed) energy states, are the differences in energy eigenvalues between this finite well and the infinite well getting ``closer and closer'' or more and more different?
          How/why?
      \end{itemize}
          \begin{answer}
            Based on the graph we drew in class, as the potential increases the energy increases asymptotically to that of the finite well.
            However, the finite well energies are all slightly smaller than the ideal infinite well.
            This can be explained by both the graph, where the asymptotic approach is from below, or by the quantum tunneling in the finite well.
            As you go higher in energy states, the differences are more different, again explained by the graph, since the finite well eventually begins to decrease down to meet the horizontal axis.
          \end{answer}
          \begin{figure}
            \includegraphics[width=\textwidth]{IMG_0351.jpeg}
          \end{figure}
  \end{enumerate}
\end{problem}

\end{document}
